This dissertation presents the formulation and simulation of a Rendezvous and Docking or Berthing mission performed by a chaser spacecraft equipped with a robotic manipulator, with emphasis on a coherent sequencing of mission phases and on consistency among models, reference frames, and control laws throughout the operational profile. The implementation was carried out in MATLAB/Simulink, using models and controllers developed by the author, without relying on ready-made RVDB simulators or predefined modules. The long-range phase is represented by a Hohmann transfer between coplanar circular orbits, used to establish initial conditions for the short-range approach; in this phase, relative motion is described by the linearized Hill--Clohessy--Wiltshire equations and regulated by a proportional--derivative control law, while the chaser attitude dynamics are modeled by Euler’s rigid-body equation. During the final approach, the dynamic model is switched to a coupled spacecraft--manipulator system derived from Lagrange’s equations, and a PD attitude controller acts to mitigate disturbances induced by joint motion; the manipulator is commanded through an inverse-kinematics reference. Performance is assessed through end-effector position error metrics and geometric alignment metrics between the tool axis and the docking port normal, both expressed in the manipulator reference frame. The results indicate that the proposed architecture enables stable and consistent phase transitions, supporting controlled approach and the geometric alignment required in the studied scenario, while highlighting the importance of model switching and disturbance decoupling/compensation in the final phase. Study limitations include the use of linearized relative dynamics for coplanar circular orbits and the absence of environmental perturbations, which restricts direct generalization to more complex operational scenarios.